\documentclass[11pt]{article}
\usepackage[T1]{fontenc}
\usepackage{lmodern,amsmath,amssymb,amsthm,mathtools,microtype}
\usepackage[hmargin=0.90in,vmargin=0.76in]{geometry}
\usepackage[hidelinks]{hyperref}
\hypersetup{pdftitle={The jump set of the one-scale problem for iid linear inequalities},pdfauthor={Peiwei Wu and Qian Peng}}
\newtheorem{theorem}{Theorem}
\newtheorem{lemma}[theorem]{Lemma}
\newtheorem{corollary}[theorem]{Corollary}
\newcommand{\bbP}{\mathbb P}
\newcommand{\supp}{\operatorname{supp}}
\newcommand{\ph}{\varphi}
\title{The jump set of the one-scale problem\\for iid linear inequalities}
\author{Peiwei Wu\qquad Qian Peng}
\date{}
\begin{document}
\maketitle
\begin{abstract}
For independent draws from a common law supported on $[a,1]$, let
$\beta(c)$ be the supremum of $\bbP(X_1+X_2<cX_3)$ over laws with
$ca>1$. We determine every jump of $\beta$ strictly between one and
the golden ratio $\ph$. The jump set consists exactly of the two
interleaved ladder families of Zaytsev and Bellec, and each jump equals
the difference of two successive finite cubic maxima. At $\ph$ the
supremum is $1/4$ and no admissible law attains it. For arbitrary
nonnegative laws we give a probability bound in terms of pairwise
scale separation, with sharp first two uniform coefficients. The
unrestricted optimum approaches $2/3$ at inverse-logarithmic rate
as $c$ increases to two.
\end{abstract}

\section{The one-scale extremum}
For independent nonnegative random variables $X_1,X_2,X_3$ of common
Borel law $\mu$, write $P_c(\mu)=\bbP(X_1+X_2<cX_3)$ and
$S_2(c)=\sup_\mu P_c(\mu)$. Zaytsev and Bellec obtained eighteen
discontinuities of $S_2$ and identified several families of candidate
thresholds \cite[Theorem~1.3 and \S\S6.1.4--6.1.7]{zb}. Their
one-scale problem is
\[
 \beta(c)=\sup\{P_c(\mu):\supp\mu\subseteq[a,1]
                       \text{ for some }a>0\text{ with }ca>1\}.
\]
A construction repeating the same distribution at widely separated
scales turns a one-scale success probability $b$ into the lower bound
$S_2(c)\ge V(b)$, where
$V(b)=\sup_{0<w<1}(1-(1-b)w^2)/(3-3w+w^2)$, for $1<c<2$
\cite[Proposition~6.1]{zb}. Whether $S_2(c)=V(\beta(c))$ is asked in
\cite[\S7.3]{zb}.

\begin{theorem}\label{th:scale}
Put $\ph=(1+\sqrt5)/2$. For each $1<C<\ph$, there exists $B(C)$ such
that, for every $1<c\le C$, a law on at most $B(C)$ atoms attains
$\beta(c)$. The function takes only finitely many values on $(1,C]$.
For every $1<a<\ph$, it has a positive right jump in $(a,\ph)$.
Consequently its jump points accumulate at $\ph$ from below.
\end{theorem}

Previous constructions introduce new winning triples at infinitely
many coefficients approaching $\ph$ \cite[\S\S6.1.4--6.1.5]{zb}.
A new winning triple in a prescribed construction yields a lower bound.
Theorem~\ref{th:scale} concerns changes in the optimum over all
one-scale laws. Theorem~\ref{th:location} identifies exactly which coefficients give jumps and expresses each jump size as a difference of two finite maxima.

Multiplying all observations by a positive constant preserves every
comparison. Dividing by $a$ permits us to work on $[1,u]$ with $u<c$.
Conversely, dividing by $u$ restores the original normalization.
For a finite configuration $1\le x_i<c$ with masses $w_i\ge0$ summing
to one, its score is
\[
  \sum_{i,j,k}w_iw_jw_k\,\mathbf1\{x_i+x_j<cx_k\}.
\]
These scores have supremum $\beta(c)$. Indeed, take an iid sample of
$N\ge3$ points from any admissible Borel law and form its equally
weighted empirical law. In expectation, the terms indexed by three
distinct sample labels contribute $(N-1)(N-2)P_c(\mu)/N^2$.
All remaining terms are nonnegative. The empirical law is admissible
almost surely, so the claim follows upon letting $N$ increase. This
empirical comparison is the sampling argument of Bellec and Fritz
\cite[Theorem~2.4 and Remark~2.5]{bf}.

\section{Compression and attainment}
\begin{lemma}\label{le:compact}
Fix $1<C<\ph$. Every finite configuration at any $1<c\le C$ can be
replaced by one on at most $B(C)$ distinct points while retaining all
its winning comparisons.
\end{lemma}
\begin{proof}
Choose $d\in(1,c)$ sufficiently close to $c$ that all the finitely
many comparisons won at $c$ still hold at $d$. Value each rooted ordered full binary tree recursively. A leaf has
value $1$ and a parent with child values $u,v$ has value $(u+v)/d$. All values are at least
$1$. Set $\delta=(1+C-C^2)/C>0$. If a parent has value at most $C$,
then its children have values at most $C^2-1<C$. Moreover
\[
 \frac{u+v}{d}-u
   =\frac{v-(d-1)u}{d}
   \ge \frac{1-(C-1)C}{C}=\delta .
\]
The analogous bound holds for the other child. Along a branch of
height $h$ in a tree of value at most $C$ the total increase is at
least $h\delta$. Only finitely many tree shapes can therefore have
value at most $C$, uniformly in $d$. Take $B(C)$ to be the number of
shapes of height at most $\lfloor(C-1)/\delta\rfloor$.

Replace $x_i$ by the largest tree value $y_i\le x_i$. The leaf is
admissible. For an old win $i,j,k$, the parent of trees giving $y_i,y_j$
has value $(y_i+y_j)/d<x_k$. Hence $y_i+y_j\le d y_k<c y_k$.
Coincident values may now be merged by adding their masses.
\end{proof}

For a fixed number $B$ of labels, each configuration determines a
subset $H$ of the $B^3$ ordered triples. For that subset the score is
the polynomial $\sum_{(i,j,k)\in H}w_iw_jw_k$ on the probability
simplex. Write $M_H$ for its maximum, which exists by compactness.
By Lemma~\ref{le:compact}, $\beta(c)$ is the largest $M_H$ among the
patterns realizable at $c$, using zero masses to pad smaller
configurations. This finite maximum is attained. The numbers $M_H$
are independent of $c$, proving the first two assertions of
Theorem~\ref{th:scale}.

\section{Location of the jumps}
For $1<c<\ph$, define
\[
 r_0(c)=1,\qquad b_0(c)=4/c^2,\qquad
 r_{n+1}(c)=\frac{1+r_n(c)}{c},\qquad
 b_{j+1}(c)=\frac{1+b_j(c)}{c}.
\]
These are the ordinary and branching ladders of Zaytsev and Bellec
\cite[\S\S6.1.4--6.1.5]{zb}. Their thresholds $c_n$ and $\tau_j$
satisfy $r_n(c_n)=c_n$ and $b_j(\tau_j)=\tau_j$, respectively.

Write $v_T(s)$ for the value of a binary tree $T$ at coefficient $s$.
We first classify the values $v_T(c)\le c$. Each child of such a root
has value at most $c^2-1<c$, since the other child has value at least $1$.
The smallest value of an internal node is $r_1=2/c$. If at least one
of its children is internal, its value is at least $r_2=(1+2/c)/c$.
But $r_1+r_2>c^2$, since
\[
 c^4-3c-2=(c^2-c-1)(c^2+c+2)<0 .
\]
Thus a node of value at most $c$ with two internal children must
have value $(r_1+r_1)/c=b_0$. Every other internal node has a leaf
child and applies $u\mapsto(1+u)/c$ to the other child's value.
Induction gives $v_T(c)\in\{r_n(c),b_j(c):n,j\ge0\}$.

The same two ladders give an exact finite optimization at each coefficient.
Write $F(u)=(1+u)/c$ and let $A_c$ consist of the values $r_n(c),b_j(c)$
strictly less than $c$, with repetitions removed.

\begin{theorem}\label{th:canonical}
For $1<c<\ph$, the set $A_c$ is finite and
\[
 \beta(c)=\max_{\substack{w_x\ge0\\\sum_{x\in A_c}w_x=1}}
   \sum_{\substack{x,y,z\in A_c\\x+y\le cz}}w_xw_yw_z .
\]
A triple contributes precisely when one of the following holds:
$x=1$ and $F(y)\le z$; $y=1$ and $F(x)\le z$; or
$x=y=r_1$ and $b_0\le z$.
\end{theorem}
\begin{proof}
Every level in $A_c$ is a tree value, so the height bound implies finiteness.
Round each atom $x_i$ of a finite law down to the largest tree value $z_i\le x_i$
at the coefficient $c$ itself. If $x_i+x_j\le cx_k$, joining the two
trees at $i,j$ shows that $(z_i+z_j)/c\le z_k$. Merging coincident
levels proves one inequality. Conversely fix the full finite set $A_c$
and apply $f_K(z)=1+K(z-1)$ with $K>1$ sufficiently close to one. All
images still lie in $[1,c)$, and
\[
 cf_K(z)-f_K(x)-f_K(y)=K(cz-x-y)+(K-1)(2-c).
\]
Each weak win becomes strict, and each strictly losing triple remains
losing by continuity. Thus the displayed cubic is realized for any
choice of weights, including its simplex maximum.
If $x,y>1$, both are at least $r_1$; unless both equal $r_1$,
one is at least $r_2$. The inequality $r_1+r_2>c^2$ obtained above
then rules out the triple. For two copies of $r_1$ the first possible
target is $2r_1/c=b_0$.
\end{proof}

\section{Every ladder threshold gives a jump}
Put the ordinary and branching levels in increasing order, indexing from
zero. Their order is
\[
 r_0,r_1,r_2,r_3,b_0,r_4,b_1,r_5,b_2,\ldots .
\]
Indeed, $(b_0-r_3)c^3=(c-1)(2-c)$ and
$(r_4-b_0)c^4=(2-c)(1+c-c^2)$ are positive, and $F$ is strictly
increasing. Denote the $i$th level by $z_i(c)$. Let $N$ be the
number strictly below $c$ and $M_N$ be the cubic maximum from
Theorem~\ref{th:canonical} on its first $N$ ranks.
Its coefficients are independent of $c$: a pair with one rank zero
and other rank $j$ wins against $k$ precisely when
$k\ge j+1$ for $j<3$, or $k\ge j+2$ for $j\ge3$.
Two inputs of rank one win against $k\ge4$.
These are all the wins.

\begin{theorem}\label{th:location}
In $(1,\ph)$ the jump set of $\beta$ is exactly
$\{c_n:n\ge1\}\cup\{\tau_j:j\ge0\}$.
When $z_N(c)=c$, the value at $c$ is $M_N$ and the right jump has
size $M_{N+1}-M_N>0$.
\end{theorem}

For $N\ge3$ the problem reduces to a cubic in four aggregate masses.
If $N=3+2k$, take $\ell=1$ and $k$ bonus slots. If $N=4+2k$,
take $\ell=2$ and $k$ bonus slots. In each case $M_N$ is the maximum of
\[
 G_{\ell,k}(p,u,v,h)
 =p(1-p)-p\left(u^2+\frac{v^2}{\ell}+\frac{h^2}{k}\right)+u^2h
\]
over nonnegative $p,u,v,h$ with sum one. For $k=0$, require $h=0$ and
omit $h^2/k$. Here $p,u$ are the bottom two masses; $v$ is the total on
the ordinary slots, and $h$ is the total on slots that win against
two rank-one inputs.

\begin{proof}[Proof of the reduction]
Write $w_i$ for the mass at rank $i$. In addition to the bonus term
$u^2\sum_{i\ge4}w_i$, the score is
\[
 p(1-p)-p\sum_{i\ge1}w_i^2-2p\sum_{i\ge3}w_iw_{i+1},
\]
where missing weights are zero. In the tail starting at rank three,
move mass between two adjacent occupied ranks. With their sum fixed,
the two square terms and their shared product have constant sum; the
remaining dependence on the displaced mass is affine. One endpoint
keeps or improves the score and deletes a tail rank. Repeating leaves
pairwise nonadjacent occupied tail ranks.

A tail of length $2k$ has at most $k$ occupied ranks. Moving their
weights onto $4,6,\ldots,2k+2$ makes all of them eligible for the
bonus. Rank two is the sole ordinary slot. A tail of length $2k+1$
has at most $k$ occupied ranks above rank three. Treat ranks two and
three as ordinary, and place the bonus weights on
$5,7,\ldots,2k+3$. In either case equalizing the weights within each
class minimizes the sum of squares. Conversely these displayed ranks,
with the ordinary slots, realize the equal weights for every admissible
$(p,u,v,h)$. This gives the asserted maximum.
\end{proof}

\begin{lemma}
At any maximum of $G_{\ell,k}$ of value at least $5/24$,
the coordinates $p,u,v$ are positive, as is $h$ whenever $k>0$.
\end{lemma}
\begin{proof}
First $p>1/4$. Otherwise put $z=1-2p\in[1/2,1]$. The objective is
at most
$(1-z^2)/4+u^2(z-u)\le(1-z^2)/4+4z^3/27$.
The last expression is decreasing on this interval and its value at
$1/2$ is $89/432<5/24$.

Let $g_p,g_u,g_v,g_h$ denote the four partial derivatives. A transfer
from a positive coordinate to another feasible coordinate cannot
increase the objective to first order. If $u=0$ and $v>0$,
$g_u=0>g_v$; if $u=v=0$ and $h>0$, $g_u=0>g_h$.
The remaining case has $p=1$ and score zero. Thus $u>0$.
If $k>0$ and $h=0$, then $g_h=u^2>g_u=-2pu$, again impossible.

Suppose $v=0$. Moving mass from $u$ to $v$ shows
$g_u=2u(h-p)\ge g_v=0$. Hence $h\ge p>1/4$, $k>0$, and $u<1/2$.
Put $q=h/k>0$. Since $p,u,h>0$, their three partial derivatives
coincide. Equating $g_u$ and $g_h$, then $g_p$ and $g_h$, gives
\[
 u^2=2u(h-p)+2pq,\qquad
 u+(h-p)=2u^2+(h-2p)q .
\]
The first gives $h-p<u/2<1/4<p$. In the second identity, the left
side is at least $u$ and the right side is strictly below $u$.
\end{proof}

\begin{proof}[Proof of Theorem~\ref{th:location}]
We show that $M_{N+1}>M_N$ for every $N\ge1$. For $N=1$ the
score vanishes. For $N=2$ it is $p^2(1-p)$, with maximum $4/27$.
For $N=3$ the masses $(1/2,1/4,1/4)$ give $3/16$, and equalizing
the two nonbottom masses gives the upper bound
$\max_{0\le p\le1}(p-p^3)/2=1/(3\sqrt3)<1/5$.
Four ranks with weights $(1/2,1/6,1/6,1/6)$ give $5/24$.
The same weights padded by zeros give $M_N\ge5/24$ for $N\ge4$.

When a new ordinary slot becomes available, replacing $v^2$ by $v^2/2$
increases the value at the old maximizer by $pv^2/2>0$.
At the other transition the reduced model loses one ordinary slot and gains one bonus slot.
Keep $p,u$ and replace $(v,h)$ by $(v/2,h+v/2)$. The penalty for the
squared slot masses cannot increase: one formerly ordinary slot
joins the $k$ bonus slots and the new group is equalized. The bonus
increases by $u^2v/2>0$. This proves strict growth.

At a root $z_N(c)=c$, there are exactly $N$ active ranks. Every
$z_i(s)$ is continuous and nonincreasing in $s$. Thus immediately
to the right of $c$, rank $N$ is active and rank $N+1$ is still
inactive. Conversely, away from roots the first inactive rank stays
inactive and the last active rank stays active in a neighborhood;
the cubic is constant there. Each $z_i$ with $i\ge1$ starts above
one at $s=1$ and is below $\ph$ at $s=\ph$, by its recurrence.
It has a unique crossing with $s$, and the ordering of the ranks
orders the crossings. The finite support bound on $(1,C]$, or the
tree height bound on all active levels, forces the crossings to
accumulate at $\ph$. This also completes Theorem~\ref{th:scale}.
\end{proof}

\begin{corollary}\label{co:golden}
At the golden ratio, $\displaystyle\lim_{c\uparrow\ph}\beta(c)=\beta(\ph)=1/4$.
No admissible Borel law attains $\beta(\ph)$. For $d>\ph$, $\beta(d)\ge3/8$.
\end{corollary}
\begin{proof}
Take $1<c\le\ph$, and split $[1,c)$ at $r_1=2/c$ and
$r_2=(1+r_1)/c$ into a lower, middle and upper band, placing
boundary points in the higher band. Write $p,t$ for the masses of
the first two bands. Every winner exceeds both inputs since
$(c-1)x<1$ for $x<c$. The lower band contains no winning targets.
With one lower input, the other input and the target cannot both
lie in the middle band, since their comparison would require
$1+r_1<c r_2$. With no lower input, both inputs must be in the
middle band and the target in the upper band: $r_1+r_2\ge c^2$
and $2r_1>c r_2$. Pairing the two possible orders of the nonlower
observations in the one-lower-input case gives
\[
 P_c(\mu)\le p^2(1-p)+p\big[(1-p)^2-t^2\big]+t^2(1-p-t)
          =p(1-p)+t^2(1-2p-t).
\]
If $u=1-2p<0$ this is at most $1/4$. For $0\le u\le 1$ use
$t^2(u-t)\le4u^3/27$ and $p(1-p)=(1-u^2)/4$ to obtain the same bound.

For the reverse bound, put half the mass at $x_0=1$ and distribute
the rest equally among $x_j=\ph-(\ph-1)2^{-j}$, $1\le j\le n$.
Here $1+x_j<\ph x_{j+1}$ for $0\le j<n$, by $\ph^2=1+\ph$
and $\ph<2$. A pair of bottom atoms wins against every higher
target. A bottom atom paired with $x_j$, $j\ge1$, wins against
every $x_k$ with $k>j$. These events contribute
$1/8+(n-1)/(8n)=1/4-1/(8n)$. Each finite collection of strict
comparisons persists at coefficients slightly below $\ph$, as
does $x_n<c$. Now let $n$ tend to infinity.

For nonattainment, fix $1\le u<\ph$. Choose $K>1$ such that
$U=1+K(u-1)<\ph$, then $\max\{1,U\}<d<\ph$ so close to $\ph$
that $(\ph-d)U<(K-1)(2-\ph)$. For $x,y,z\in[1,u]$,
\[
 d f_K(z)-f_K(x)-f_K(y)
 =K(\ph z-x-y)+(K-1)(2-\ph)-(\ph-d)f_K(z).
\]
Every win at $\ph$ persists at $d$ after the change of variables.
Thus $P_{\ph}(\mu)\le\beta(d)$ uniformly for laws on $[1,u]$.
There is a positive jump between $d$ and $\ph$, by
Theorem~\ref{th:scale}, so $\beta(d)<1/4$. Finally, for $d>\ph$
choose $\max\{1,1/(d-1)\}<T<d$. With equal weights on $1,T$, at least the
three triples $(1,1,T),(1,T,T),(T,1,T)$ win among eight equally likely
ordered triples, as in \cite[Corollary~6.2]{zb}.
\end{proof}

\begin{corollary}\label{co:global}
For the unrestricted Borel supremum,
$\displaystyle S_2(\ph)\ge (\sqrt{61}-5)/6$.
\end{corollary}
\begin{proof}
We use the separated-scale construction of \cite[Proposition~6.1]{zb}.
Let $t$ be the upper block mass and $p=1-t$, where $0<t<1$.
Place any finite one-scale law on $[1,u]$, $u<c$, in the upper block.
Rescale a finite nonnegative law into $[0,\varepsilon]$ for the lower
block, with $2\varepsilon<c$ and $\varepsilon+u<c$.
When the target is in the upper block, two lower inputs always win,
as does one lower input paired with an upper input. Keeping just
these wins and those internal to each block, then taking suprema,
gives
\[
 (1-p^3)S_2(c)\ge t^3\beta(c)+2pt^2+p^2t .
\]
Finite laws approximate both suprema by the empirical comparison
of Section~1. At $c=\ph$, insert $\beta(\ph)=1/4$ and take
$t=(13-\sqrt{61})/9$. Substitution gives the asserted bound.
\end{proof}
The numerical consequence $S_2(\ph)\ge 0.46837494$ also follows directly by taking increasingly long ordinary ladders in \cite[\S6.1.4]{zb} and applying their Proposition~6.1. For comparison, the
six-point endpoint certificate of Zaytsev and Bellec gives
$S_2(\ph)\le 846522430902936932843/(18\cdot10^{20})<0.470291$
\cite[\S6.1.8 and ancillary certificates]{zb}.

\section{A bound for laws occupying several scales}
For independent $X,Z$ of common law $\mu$, put
$D_c(\mu)=\bbP(cX\le Z)$. This counts ordered pairs whose larger
entry lies outside the admissible one-scale range of the smaller entry.

\begin{theorem}\label{th:spread}
Let $1<c\le\ph$. For every nonnegative Borel law $\mu$,
\[
 P_c(\mu)\le\beta(c)+2D_c(\mu)-D_c(\mu)^2 .
\]
\end{theorem}
\begin{proof}
First take a finite law. Call a pair $a,b$ compatible if $a<cb$ and
$b<ca$. Retain the wins whose three pairs are all compatible, and
denote their total weight by $C$. If two positive-mass labels are
incompatible, no retained monomial contains both. On transferring
mass between them with the sum of their masses fixed, each monomial
is a nonnegative multiple of a constant, a power of the first mass,
or a power of the second. Thus $C$ is convex on the transfer interval;
one endpoint preserves or increases $C$ and removes a label.
Iterate. A remaining support with at least two labels is strictly
positive and its extreme values $a,b$ satisfy $b<ca$. Rescaling by
$a$ now gives a one-scale law, so $C\le\beta(c)$. A singleton at
zero contributes no wins.

In a winning triple $x+y<cz$, the target is positive. If $cx>z$ and
$cy>z$, then $x,y>z/c$ and $x,y<(c-1/c)z\le z$. All three pairs
are then compatible. Every win omitted from $C$ therefore has
$cx\le z$ or $cy\le z$. Put $q(z)=\bbP(cX\le z)$. Conditional on
$z$, the probability of this union is $2q(z)-q(z)^2$. Averaging and
using $\mathbb E q(Z)^2\ge (\mathbb E q(Z))^2$ proves the bound.

For a Borel law let $\widehat\mu_N$ be the empirical law of $N$
independent draws. The expected empirical triple score is at least
$N(N-1)(N-2)P_c(\mu)/N^3$, by retaining distinct indices.
The expected empirical pair probability tends to $D_c(\mu)$:
off the diagonal it equals $D_c(\mu)$, and diagonal pairs have total
weight $1/N$. Apply the finite inequality to $\widehat\mu_N$, take
expectations and use concavity of $d\mapsto2d-d^2$. Sending
$N$ to infinity gives the claim.
\end{proof}

The linear and quadratic coefficients are sharp for a bound
of the form $\beta(c)+AD-BD^2$ uniform in $1<c\le\ph$.
At $c=4/3$ no one-scale triple wins, since $c^2<2$. Put probabilities
$p,1-p$ on $1,4$. Then $D=p(1-p)$ and $P_c=(2-p)D$.
Letting $p\downarrow0$ after division by $D$ forces $A\ge2$; when
$A=2$, the same inequality gives $B(1-p)\le1$ and hence $B\le1$.

\section{The unrestricted problem near two}
The arithmetic restrictions on locations of unrestricted jumps
\cite{wp} do not determine where the supremum varies. The next estimate
gives its rate of approach to the endpoint.
\begin{theorem}\label{th:end}
Let $1<c<2$ and $L=\log(1/(2-c))/\log c$. Then
\[
 \frac{1}{3(L+1)}
 \ \le\ \frac23-S_2(c)
 \ \le\ \frac{1}{2L}+\frac{1}{6L^2}.
\]
In particular $2/3-S_2(c)=\Theta(1/\log(1/(2-c)))$ as $c\uparrow2$.
\end{theorem}
\begin{proof}
Put $m=\lceil L\rceil$, so $c^{-m}\le2-c<c^{-(m-1)}$. Write
$t_0=0$ and $t_{j+1}=(1+t_j)/c$. The identity
$t_m=(1-c^{-m})/(c-1)$ gives $t_m\ge1$.

First consider an atomless law. Order three draws and condition on their
maximum $z$. The two smaller observations in random order are
independent draws from the law truncated at $z$.
Partition $[0,z]$ at $t_1z,\ldots,t_{m-1}z$. If $x\le y$ lie in
the same interval, say the $j$th, then
$c y\le c t_{j+1}z=(1+t_j)z\le z+x$. The middle observation then
cannot win. The smallest observation can never win, since $c<2$.
The conditional probability of both lower draws falling in the same
class is at least $1/m$ by Cauchy--Schwarz. Exchangeability gives
$3P_c(\mu)\le 2-1/m$. For an arbitrary law add independent uniform
noise on $[0,\varepsilon]$ to each draw. The perturbed law is atomless,
and every strict win persists for all sufficiently small
$\varepsilon$. Fatou's lemma yields the same bound.

For the other direction use the equal-weight ladder with $m$ levels
$x_j=1-c^{-j}$ for $0\le j<m$. This is the pattern of
\cite[\S6.1.7]{zb}. We have $x_{m-1}<c-1$ and
$cx_{j+1}=x_j+c-1$. A comparison $x_i+x_j<cx_k$ therefore holds
exactly when $\min(i,j)<k$. Against target $k$, precisely $(m-k)^2$
ordered pairs lose. Summing over $k$ gives
\[
 P_c(\mu)=1-\frac{1}{m^3}\sum_{h=1}^m h^2
          =\frac23-\frac{1}{2m}-\frac{1}{6m^2}.
\]
The same polynomial in $1/m$ appears in the endpoint game of
Bellec and Fritz \cite[Theorem~1.5]{bf}. Combining these bounds with
$L\le m<L+1$ proves the theorem.
\end{proof}

Let $F_N(c)$ be the optimum for a uniform law on $N$ labelled points,
allowing repetitions. The unrestricted problem and these finite
profiles obey
\[
  F_N(c)\le S_2(c)\le F_N(c)+1-\frac{N(N-1)(N-2)}{N^3}
  \qquad(N\ge3).
\]
Theorem~\ref{th:end} implies that $S_2$ takes infinitely many distinct
values in each left neighborhood of $2$. In particular, given $a<2$ and
a positive integer $K$, one can choose $K$ disjoint intervals
$[s_i,t_i)\subset(a,2)$ with $S_2(t_i)>S_2(s_i)$. For all sufficiently
large $N$ the same strict inequalities hold for $F_N$. Its values are
multiples of $N^{-3}$ and its maximum at each coefficient persists just
to the left, by strictness of the winning comparisons in a maximizing
sample. Hence each of these intervals contains a right jump of $F_N$
of size at least $N^{-3}$. For $S_2$ itself, whether there are
infinitely many jumps remains the question of \cite[\S7.1.3]{zb}.

\begin{thebibliography}{9}
\bibitem{bf} P.~C. Bellec and T. Fritz,
On maximizing the probability of a linear inequality between iid random variables,
\emph{Amer. Math. Monthly} (2026),
\href{https://doi.org/10.1080/00029890.2026.2702871}{doi:10.1080/00029890.2026.2702871}.
\bibitem{wp} P. Wu and Q. Peng, Spectral restrictions on the jumps of extremal iid probabilities,
preprint (2026), \href{https://github.com/PEIWEIWU-AYS/iid-extremal-probability-jumps-spectral-radii/releases/tag/v1.0.0}{version 1.0.0}.
\bibitem{zb} A. Zaytsev and P.~C. Bellec,
Flag algebras and the exact largest probability of a linear inequality between iid random variables,
\href{https://arxiv.org/abs/2611.29223}{arXiv:2611.29223} (2026).
\end{thebibliography}
\end{document}
